% from ega0-19b-fr.tex
\oldpage[0\textsubscript{IV-1}]{115}

\begin{env}[19.9.5]
\label{0.19.9.5-fr}
De même, l'énoncé (19.3.8) est encore valable (avec la même démonstration) quand
on y remplace «~formellement lisse~» par «~formellement lisse relativement à
$\Lambda$~». Si dans l'énoncé de (19.3.10) on remplace «~formellement lisse~»
par «~formellement lisse relativement à $\Lambda$~» la conclusion est remplacée
par la suivante (la démonstration restant essentiellement inchangée)~: \emph{tout
$A$-homomorphisme continu $u_0:B\to C/\mathfrak J$ qui se factorise en
$B\xrightarrow{u}C\to C/\mathfrak J$, où $u$ est un $\Lambda$-homomorphisme
continu, se factorise aussi en $B\xrightarrow{v}C\to C/\mathfrak J$, où $v$
est un $A$-homomorphisme continu}.
\end{env}

\begin{env}[19.9.6]
\label{0.19.9.6-fr}
Les critères de lissité formelle (19.4.1) et (19.4.2) sont valables lorsqu'on y
remplace «~formellement lisse~» par «~formellement lisse relativement à
$\Lambda$~», les démonstrations restant pratiquement inchangées.
\end{env}

\begin{proposition}[19.9.7]
\label{0.19.9.7-fr}
Soient $\Lambda$ un anneau topologique, $A$ une $\Lambda$-algèbre topologique,
$B$ une $A$-algèbre topologique. Supposons que pour toute $A$-algèbre discrète
$C$ et tout idéal $\mathfrak J$ de $C$ tel que $\mathfrak J^2=0$, tout
$A$-homomorphisme continu $u_0:B\to C/\mathfrak J$ qui se factorise en
$B\xrightarrow{u}C\to C/\mathfrak J$, où $u$ est un $\Lambda$-homomorphisme
continu, se factorise aussi en $B\xrightarrow{v}C\to C/\mathfrak J$, où $v$
est un $A$-homomorphisme continu. Alors $B$ est une $A$-algèbre formellement
lisse relativement à $\Lambda$.
\end{proposition}

\begin{proof}
La démonstration de (19.4.3) se transcrit aussitôt.
\end{proof}

\begin{proposition}[19.9.8]
\label{0.19.9.8-fr}
Soient $\Lambda$ un anneau topologique, $A$ une $\Lambda$-algèbre topologique,
$B$ une $A$-algèbre topologique. Pour que $B$ soit une $A$-algèbre
formellement lisse relativement à $\Lambda$, il faut et il suffit que pour tout
$B$-module topologique discret $L$, annulé par un idéal ouvert de $B$, on ait
(cf.~18.4.2)
\[
  \operatorname{Exalcotop}_{A/\Lambda}(B,L)=0.
\]
\end{proposition}

\begin{proof}
Avec les notations de la démonstration de (19.4.4), il suffit de noter ici que
l'on peut supposer que l'extension $E_\Lambda$ de $B_\Lambda$ est
$\Lambda$-triviale~; le reste de la démonstration est alors inchangé.
\end{proof}

Lorsque $\Lambda$, $A$ et $B$ sont des anneaux discrets, le critère (19.9.8) se
réduit à
\begin{equation}
\tag{19.9.8.1}
\label{0.19.9.8.1-fr}
  \operatorname{Exalcom}_{A/\Lambda}(B,L)=0
  \quad\text{pour tout $B$-module $L$}~;
\end{equation}
autrement dit, \emph{toute $A$-extension commutative de $B$ par un $B$-module,
qui est $\Lambda$-triviale, est aussi $A$-triviale}.

\subsection{Algèbres formellement non ramifiées et algèbres formellement
étales}
\label{subsection:0.19.10-fr}
\label{0.19.10-fr}

\begin{env}[19.10.1]
\label{0.19.10.1-fr}
Nous aurons aussi à utiliser deux notions voisines de la notion d'algèbre
formellement lisse.
\end{env}

\begin{definition}[19.10.2]
\label{0.19.10.2-fr}
Soient $A$ un anneau topologique, $B$ une $A$-algèbre topologique. On dit que
$B$ est une $A$-algèbre \emph{formellement non ramifiée} si, pour tout
$A$-homomorphisme $p:E\to C$ de $A$-algèbres discrètes, dont le noyau est
nilpotent, et tout $A$-homomorphisme continu $u:B\to C$, il existe \emph{au plus
un} $A$-homomorphisme continu $v:B\to E$ tel que $u$ se factorise en
$B\xrightarrow{v}E\xrightarrow{p}C$. On dit que $B$ est une $A$-algèbre
\emph{formellement étale} si elle est formellement lisse et formellement non
ramifiée, autrement dit, si, pour tout $A$-homomorphisme surjectif $p:E\to C$
de $A$-algèbres discrètes, dont le noyau est nilpotent, et pour tout
$A$-homomorphisme continu $u:B\to C$, il existe \emph{un et un seul}
$A$-homomorphisme continu $v:B\to E$ tel que $u$ se factorise en
$B\xrightarrow{v}E\xrightarrow{p}C$.
\end{definition}

